\chapter*{Glossary and abbreviations}
\addcontentsline{toc}{chapter}{Glossary and abbreviations}

\begin{description}
\item[AI] Adequate Intake, a reference intake used when an RDA cannot be established.
\item[DFE] Dietary folate equivalents, a unit accounting for differences between food folate and folic acid.
\item[DRI] Dietary Reference Intake, a family of reference values including EAR, RDA, AI, and UL.
\item[EAR] Estimated Average Requirement for a defined population group.
\item[NE] Niacin equivalents.
\item[RAE] Retinol activity equivalents.
\item[RDA] Recommended Dietary Allowance, intended to meet the needs of nearly all healthy people in a group.
\item[UL] Tolerable Upper Intake Level, the highest average daily intake unlikely to pose risk for most people.
\item[25(OH)D] Serum 25-hydroxyvitamin D, the usual vitamin D status marker.
\item[FMN/FAD] Flavin mononucleotide and flavin adenine dinucleotide, active riboflavin-derived cofactors.
\item[NAD/NADP] Nicotinamide adenine dinucleotide and its phosphate, niacin-derived redox cofactors.
\end{description}

This glossary is educational. Laboratory units, reference intervals, and treatment thresholds must be interpreted using the local laboratory and current clinical guideline.
